% Document/LinearAlgebra/VectorSpaces/Foundations/Flags.tex
% Section: Flags and Chains of Subspaces

\newpage
\section{Flags and Chains of Subspaces}\label{sec:flags}

\begin{remark}[Motivation]
    Having introduced ordered bases, we now study a related but distinct concept: ordered
    families of \emph{subspaces}. Just as an ordered basis imposes a total order on vectors
    to produce coordinates, a \emph{flag} imposes a total order on a chain of subspaces to
    capture the geometry of nested substructures. Flags arise naturally in the study of
    triangular and block-triangular matrix forms, Gaussian elimination, and Jordan normal forms.
\end{remark}

% ─────────────────────────────────────────────────────────────────────────────
\subsection{Ordered Bases and Set-Theoretic Requirements}

\begin{remark}[Total Order vs.\ Well-Order]
    A \emph{total order} on the basis already suffices for coordinates, and \emph{no} order at
    all is needed to define matrices or matrix multiplication --- only column-finiteness (see
    the chapter on Matrices). A \emph{well-order} is needed only for the transfinite
    \emph{constructions} of this section (recovering an adapted basis from a complete flag) and
    of later chapters (echelon and rank normal forms, Gram--Schmidt). For
    \textbf{finite-dimensional} spaces every total order is a well-order, so the distinction
    vanishes. See Remark~\ref{rem:order-buys} for the full accounting of what each layer of the
    theory requires.
\end{remark}

% ─────────────────────────────────────────────────────────────────────────────
\subsection{Flags}

\begin{definition}[Flag]\label{def:flag}
    A \textbf{flag} in a $K$-vector space $V$ is a totally ordered set (chain) of subspaces
    ordered by strict inclusion. A \textbf{finite flag} of \textbf{length} $r+1$ is a sequence
    \[
        V_0 \subsetneq V_1 \subsetneq \cdots \subsetneq V_r
    \]
    of $r+1$ subspaces of $V$. For an \textbf{infinite flag}, one has a chain
    $(V_\alpha)_{\alpha \in A}$ indexed by a totally ordered set $A$, with $V_\alpha \subsetneq V_\beta$
    whenever $\alpha < \beta$.

    Equivalently, a flag is a chain in the poset $\bigl(\SubspaceLattice{V},\, \Subspace[K]\bigr)$ of all
    subspaces of $V$ ordered by inclusion.
\end{definition}

\begin{example}[Flags in $\bb{R}^2$]\label{ex:flag-R2}
    In $V = \bb{R}^2$ (dimension $2$), examples of flags:
    \begin{itemize}
        \item $\{0\} \subsetneq \bb{R}^2$ \quad --- length $2$, a chain of two subspaces.
        \item $\{0\} \subsetneq \Span[K]{\CanonicalVector{0}} \subsetneq \bb{R}^2$ \quad --- length $3$,
              a chain of three subspaces where $\Span[K]{\CanonicalVector{0}}$ is a line.
        \item $\Span[K]{\CanonicalVector{1}} \subsetneq \bb{R}^2$ \quad --- length $2$; a flag need not
              start at $\{0\}$ or end at $V$.
    \end{itemize}
    There are infinitely many flags of length $3$ in $\bb{R}^2$: one for each line through
    the origin.
\end{example}

\begin{example}[Flags in Polynomial Spaces]\label{ex:flag-poly}
    Let $V = K[X]_{\leqslant 2}$, the space of polynomials of degree at most $2$ (dimension $3$).
    Flags of different lengths:
    \begin{itemize}
        \item Length $2$: $\{0\} \subsetneq K[X]_{\leqslant 2}$
              (jumps directly from $\{0\}$ to $V$).
        \item Length $3$: $\{0\} \subsetneq \Span[K]{1} \subsetneq K[X]_{\leqslant 2}$
              (dimensions $0, 1, 3$; a gap of $2$ in the last step).
        \item Length $4$: $\{0\} \subsetneq \Span[K]{1} \subsetneq K[X]_{\leqslant 1}
              \subsetneq K[X]_{\leqslant 2}$
              (dimensions $0, 1, 2, 3$; each step increases by $1$).
    \end{itemize}
\end{example}

\begin{definition}[Complete Flag]\label{def:complete-flag}
    A flag $\mathcal{F}$ in $V$ is \textbf{complete} if it is a \emph{maximal chain} in the
    poset of all subspaces of $V$: no subspace can be inserted into $\mathcal{F}$ at any
    position without violating strict inclusion.

    In infinite dimensions a complete flag is, in addition, \textbf{continuous at limits}: being
    a maximal chain, it is closed under the unions and intersections of its subchains. (Indeed
    the union of any subchain is a subspace comparable to every member, so maximality returns it
    to the chain; dually for intersections.) Consequently each \emph{limit} member --- one with
    no immediate predecessor --- equals the union of its predecessors, and dually each member with
    no immediate successor equals the intersection of its successors; a \emph{jump} member instead
    exceeds its predecessor by a single dimension. (In finite dimensions there are no limit
    members and this is automatic.) This is the order-theoretic content of a \emph{complete nest}
    in operator theory.
\end{definition}

\begin{example}[Complete and Non-Complete Flags in $\bb{R}^3$]\label{ex:flag-complete-R3}
    In $V = \bb{R}^3$ (dimension $3$), consider flags from $\{0\}$ to $\bb{R}^3$:
    \begin{itemize}
        \item \textbf{Non-complete, length $2$}:
              $\{0\} \subsetneq \bb{R}^3$.
              Any line or plane can be inserted.

        \item \textbf{Non-complete, length $3$ (dimension gap)}:
              $\{0\} \subsetneq \Span[K]{\CanonicalVector{0}} \subsetneq \bb{R}^3$
              (dimensions $0, 1, 3$). Any plane $P$ with
              $\Span[K]{\CanonicalVector{0}} \subsetneq P \subsetneq \bb{R}^3$
              (e.g.\ $P = \Span[K]{\CanonicalVector{0}, \CanonicalVector{1}}$) can be inserted between the last two terms.
              Not maximal.

        \item \textbf{Complete, length $4$}:
              $\{0\} \subsetneq \Span[K]{\CanonicalVector{0}} \subsetneq \Span[K]{\CanonicalVector{0}, \CanonicalVector{1}} \subsetneq \bb{R}^3$
              (dimensions $0, 1, 2, 3$; each step $+1$; nothing can be inserted).

        \item \textbf{Another complete flag, length $4$}:
              $\{0\} \subsetneq \Span[K]{\CanonicalVector{0} + \CanonicalVector{1}} \subsetneq
               \Span[K]{\CanonicalVector{0} + \CanonicalVector{1},\, \CanonicalVector{2}} \subsetneq \bb{R}^3$.
    \end{itemize}
\end{example}

\begin{proposition}[Complete Flags in Finite Dimensions]\label{prop:complete-flag-finite}
    Let $\Dim[K]{V} = n \in \bb{N}$. A flag
    $V_0 \subsetneq V_1 \subsetneq \cdots \subsetneq V_r$
    with $V_0 = \{0\}$ and $V_r = V$
    is complete if and only if $r+1 = n+1$ (i.e.\ $r = n$) and $\Dim[K]{V_i} = i$ for all
    $i = 0, 1, \ldots, n$.
    In particular, every step increases dimension by exactly $1$.
\end{proposition}

\begin{proof}
    ($\Rightarrow$)
    Suppose the flag is complete. Since each inclusion is strict, $\Dim[K]{V_{i+1}} > \Dim[K]{V_i}$.
    If $\Dim[K]{V_{i+1}} \geq \Dim[K]{V_i} + 2$ for some $i$, pick any subspace $W$ with
    $V_i \subsetneq W \subsetneq V_{i+1}$ and $\Dim[K]{W} = \Dim[K]{V_i} + 1$
    (such $W$ exists by the basis extension theorem). Inserting $W$ contradicts maximality.
    Hence $\Dim[K]{V_{i+1}} = \Dim[K]{V_i} + 1$ for all $i$, giving $r = n$ and $\Dim[K]{V_i} = i$.

    ($\Leftarrow$)
    Suppose $r = n$ and $\Dim[K]{V_i} = i$. If $W$ satisfies $V_i \subsetneq W \Subspace[K] V_{i+1}$,
    then $i < \Dim[K]{W} \leqslant i+1$, so $\Dim[K]{W} = i+1 = \Dim[K]{V_{i+1}}$,
    hence $W = V_{i+1}$ by the equal-dimension corollary. No intermediate subspace fits.
\end{proof}

\begin{example}[Verifying Completeness in $\bb{R}^4$]\label{ex:flag-verify}
    Two flags from $\{0\}$ to $\bb{R}^4$:
    \begin{itemize}
        \item $\mathcal{F}_1$: $\{0\} \subsetneq \Span[K]{\CanonicalVector{0}} \subsetneq \Span[K]{\CanonicalVector{0},\CanonicalVector{1}}
              \subsetneq \Span[K]{\CanonicalVector{0},\CanonicalVector{1},\CanonicalVector{2}} \subsetneq \bb{R}^4$.
              Dimensions: $0,1,2,3,4$. Length $= 5 = 4+1$. Each step $+1$.
              \textbf{Complete.}

        \item $\mathcal{F}_2$: $\{0\} \subsetneq \Span[K]{\CanonicalVector{0}} \subsetneq
              \Span[K]{\CanonicalVector{0},\CanonicalVector{1},\CanonicalVector{2}} \subsetneq \bb{R}^4$.
              Dimensions: $0,1,3,4$. Length $= 4$. Step $1\to 3$ jumps by $2$:
              $\Span[K]{\CanonicalVector{0}, \CanonicalVector{1}}$ can be inserted. \textbf{Not complete.}
    \end{itemize}
\end{example}

% ─────────────────────────────────────────────────────────────────────────────
\subsection{Flags from Ordered Bases}

\begin{remark}[Why the Naive Successor Chain Fails]\label{rem:naive-flag-fails}
    One is tempted to set $V_\alpha \defeq \Span[K]{\Set{b_i \mid i \leqslant \alpha}}$ for each
    $\alpha \in I$ and string these into a chain through successors $\alpha^+$. This works in
    finite dimensions but is \textbf{incomplete in infinite dimensions}: it presupposes that
    every index has a successor and that $I$ has a minimum, and its maximality check only closes
    the \emph{successor} gaps $V_\alpha \subsetneq V_{\alpha^+}$. At a \emph{limit} index
    $\lambda$ the open segment $\Span[K]{\Set{b_i \mid i < \lambda}}$ lies strictly between all
    earlier $V_\alpha$ and $V_\lambda$ and is insertable, so the chain is \emph{not maximal}.
    This already fails for a basis of order type $\omega + 1$, and for every uncountable
    dimension. The cure is to index the flag by \emph{initial segments} rather than by elements,
    which is uniform over any total order.
\end{remark}

\begin{proposition}[Every Ordered Basis Induces a Complete Flag; Any Total Order]\label{prop:basis-flag}
    Let $\mathcal{B} = (b_i)_{i \in I}$ be an ordered basis of $V$, with $I$ totally ordered. For
    each \textbf{initial segment} $L \subseteq I$ --- a downward-closed subset, i.e.\ $i \in L$
    and $j < i$ imply $j \in L$ --- set
    \[
        W_L \defeq \Span[K]{\Set{b_i \mid i \in L}}.
    \]
    Then $\Set{W_L \mid L \text{ an initial segment of } I}$ is a complete flag of $V$.
\end{proposition}

\begin{proof}
    \textbf{Order-embedding.}
    If $L \subsetneq L'$ are initial segments, pick $i \in L' \setminus L$; then $b_i \in W_{L'}$,
    while $b_i \notin W_L$ by linear independence (every element of $W_L$ is a finite combination
    of basis vectors indexed in $L$, and $i \notin L$). Hence $W_L \subsetneq W_{L'}$. The initial
    segments of $I$, ordered by $\subseteq$, form a complete chain $\widehat{I}$ --- the cut
    completion of $I$ --- and $L \mapsto W_L$ embeds it into $\SubspaceLattice{V}$ as a chain. The
    least initial segment $L = \emptyset$ gives $W_\emptyset = \{0\}$, and the greatest $L = I$
    gives $W_I = V$.

    \textbf{Continuity.}
    Call $L$ a \emph{limit cut} if $L = \BigUnion[L'' \subsetneq L] L''$, the union of the
    strictly smaller initial segments. For such $L$,
    \[
        W_L = \Span[K]{\Set{b_i \mid i \in L}} = \BigUnion[L'' \subsetneq L] W_{L''},
    \]
    because any vector of $W_L$ uses only finitely many basis indices: each lies in some strictly
    smaller initial segment, and as those segments form a chain the finitely many of them all lie
    in their largest, a single $L'' \subsetneq L$. Thus the chain is closed under the unions of its
    subchains.

    \textbf{Maximality.}
    Let $X \Subspace[K] V$ be comparable (under $\Subspace[K]$) to every $W_L$. Put
    \[
        L^{*} \defeq \BigUnion \Set{L \mid W_L \Subspace[K] X},
    \]
    a union of initial segments, hence itself an initial segment; by continuity
    $W_{L^{*}} = \BigUnion \Set{W_L \mid W_L \Subspace[K] X} \Subspace[K] X$. For every initial
    segment $L \supsetneq L^{*}$ we have $W_L \not\Subspace[K] X$ (otherwise $L$ would be one of
    the segments unioned into $L^{*}$), so comparability forces $X \Subspace[K] W_L$; hence
    $X \Subspace[K] \BigIntersect[L \supsetneq L^{*}] W_L$. Two cases:
    \begin{itemize}
        \item If $I \setminus L^{*}$ has a least element $i$, then $L^{*} \Union \Set{i}$ is the
              smallest initial segment properly containing $L^{*}$, so
              \[
                  \BigIntersect[L \supsetneq L^{*}] W_L = W_{L^{*} \Union \Set{i}}
                  = W_{L^{*}} \DirectSum \Span[K]{b_i},
              \]
              a one-dimensional extension of $W_{L^{*}}$. From
              $W_{L^{*}} \Subspace[K] X \Subspace[K] W_{L^{*}} \DirectSum \Span[K]{b_i}$ we get
              $X \in \Set{W_{L^{*}},\ W_{L^{*}} \DirectSum \Span[K]{b_i}}$ --- in either case some
              $W_L$.
        \item Otherwise $I \setminus L^{*}$ has no least element. Suppose
              $v \in \BigIntersect[L \supsetneq L^{*}] W_L$ but $v \notin W_{L^{*}}$; then the
              finite support of $v$ meets $I \setminus L^{*}$ in a nonempty set, with least
              element $m$. As $I \setminus L^{*}$ has no least element, some
              $m' \in I \setminus L^{*}$ satisfies $m' < m$, so $L \defeq \Set{x \in I \mid x < m}$
              is an initial segment with $L^{*} \subsetneq L$ and $m \notin L$. Then
              $v \in W_L = \Span[K]{\Set{b_x \mid x < m}}$ forces the coefficient of $b_m$ in $v$
              to vanish, contradicting $m$ being in the support of $v$. Hence
              $\BigIntersect[L \supsetneq L^{*}] W_L = W_{L^{*}}$ and $X = W_{L^{*}}$.
    \end{itemize}
    Either way $X$ is already a member of the chain, so nothing can be inserted: the flag is
    maximal, hence complete.
\end{proof}

\begin{remark}[The Jumps Recover the Basis]\label{rem:jumps-recover}
    The covering pairs of $\widehat{I}$ are exactly $\Set{j \mid j < i} \prec \Set{j \mid j \leqslant i}$,
    one for each $i \in I$, and the corresponding jump
    $W_{\Set{j < i}} \subsetneq W_{\Set{j \leqslant i}}$ has one-dimensional quotient spanned by
    $b_i$. So the flag remembers $\mathcal{B}$: the basis vector $b_i$ sits in the $i$-th jump.
\end{remark}

\begin{remark}[Well-Ordered Specialization]\label{rem:wo-specialization}
    When $I$ is well-ordered (order-isomorphic to an ordinal $\kappa$), the initial segments are
    exactly the sets $\Set{i \mid i < \beta}$ for $\beta \leqslant \kappa$, and the cut-flag
    becomes
    \[
        W_\beta = \Span[K]{\Set{b_i \mid i < \beta}}, \qquad W_0 = \{0\}, \quad W_\kappa = V,
    \]
    \[
        W_{\beta+1} = W_\beta \DirectSum \Span[K]{b_\beta}, \qquad
        W_\lambda = \BigUnion[\beta < \lambda] W_\beta \quad (\lambda \text{ a limit ordinal}).
    \]
    This is the corrected form of the naive statement (Remark~\ref{rem:naive-flag-fails}): use
    \emph{strict} initial segments $\Set{i \mid i < \beta}$ and include the limit clause
    $W_\lambda = \BigUnion[\beta < \lambda] W_\beta$ in the completeness proof.
\end{remark}

\begin{example}[Induced Flags]\label{ex:induced-flag}
    The following ordered bases yield the following complete flags:
    \begin{itemize}
        \item Standard basis $(\CanonicalVector{0}, \CanonicalVector{1}, \CanonicalVector{2})$ of $\bb{R}^3$:
              $\{0\} \subsetneq \Span[K]{\CanonicalVector{0}} \subsetneq \Span[K]{\CanonicalVector{0},\CanonicalVector{1}} \subsetneq \bb{R}^3$.
        \item Rotated basis $(\CanonicalVector{0}+\CanonicalVector{1},\, \CanonicalVector{1},\, \CanonicalVector{2})$ of $\bb{R}^3$:
              $\{0\} \subsetneq \Span[K]{\CanonicalVector{0}+\CanonicalVector{1}} \subsetneq \Span[K]{\CanonicalVector{0}+\CanonicalVector{1},\, \CanonicalVector{1}}
               \subsetneq \bb{R}^3$.
        \item Canonical basis $(\CanonicalVector{n})_{n \in \bb{N}}$ of $K^{(\bb{N})}$:
              $\{0\} \subsetneq \Span[K]{\CanonicalVector{0}} \subsetneq \Span[K]{\CanonicalVector{0},\CanonicalVector{1}}
               \subsetneq \cdots$ (an infinite complete flag).
    \end{itemize}
\end{example}

\begin{example}[Cut-Flags in Infinite Dimensions]\label{ex:cut-flag-infinite}
    The cut construction specialises cleanly:
    \begin{itemize}
        \item \textbf{Standard basis of $K^n$}: the initial segments are
              $\emptyset \subsetneq \Set{0} \subsetneq \Set{0,1} \subsetneq \cdots \subsetneq
              \Set{0, \ldots, n-1}$, giving the standard complete flag of dimensions
              $0, 1, \ldots, n$.
        \item \textbf{Canonical basis of $K^{(\bb{N})}$}:
              $\{0\} \subsetneq \Span[K]{\CanonicalVector{0}} \subsetneq
              \Span[K]{\CanonicalVector{0}, \CanonicalVector{1}} \subsetneq \cdots$, with the whole
              space reached as the union $\BigUnion[n \in \bb{N}] W_{\Set{0, \ldots, n-1}}$ at the
              top. There is \emph{no top jump}: the top member is a limit, not a successor.
        \item \textbf{A $\bb{Q}$-indexed basis} $(b_q)_{q \in \bb{Q}}$: the initial segments of
              $\bb{Q}$ are its Dedekind cuts, so the flag is indexed by
              $\widehat{\bb{Q}} \approx \bb{R}$ (with endpoints), a continuum, even though the
              basis itself is countable. The jumps sit back at the rationals; the irrational cuts
              are continuous --- gap-filling, with no jump and no basis vector.
    \end{itemize}
\end{example}

% ─────────────────────────────────────────────────────────────────────────────
\subsection{From a Complete Flag Back to a Basis: Jumps}

\begin{definition}[Jump]\label{def:jump}
    In a complete flag, a \textbf{jump} is a covering pair $W^{-} \prec W$ (members with
    $W^{-} \subsetneq W$ and nothing strictly between). Maximality forces
    $\Dim[K]{\QuotientSpace{W}{W^{-}}} = 1$. A \textbf{jump-vector} is a choice of one vector in
    $W \setminus W^{-}$ for each jump.
\end{definition}

\begin{proposition}[Jump-Vectors Are Linearly Independent]\label{prop:jump-independent}
    The jump-vectors of any complete flag are linearly independent.
\end{proposition}

\begin{proof}
    Suppose $\sum_{k \in F} c_k v_k = 0$ with $F$ finite, every $c_k \neq 0$, and the $v_k$
    distinct jump-vectors at jumps $W_k^{-} \prec W_k$. The members $W_k$ are distinct (distinct
    jumps have distinct top members); among the finitely many of them, let $W_m$ be the highest in
    the chain, so each other $W_k$ is strictly below $W_m$. As the members form a chain, $W_k$ is
    comparable to $W_m^{-}$, and being strictly below $W_m$ with $W_m^{-} \prec W_m$ a covering
    pair forces $W_k \Subspace[K] W_m^{-}$; hence $v_k \in W_m^{-}$ for $k \neq m$. Then
    \[
        v_m = -\,c_m^{-1} \sum_{k \in F \setminus \Set{m}} c_k v_k \in W_m^{-},
    \]
    contradicting $v_m \in W_m \setminus W_m^{-}$.
\end{proof}

\begin{proposition}[When the Jumps Form a Basis]\label{prop:jump-recovery}
    The jump-vectors of a complete flag span $V$ --- and hence, with
    Proposition~\ref{prop:jump-independent}, form a basis \emph{adapted} to the flag --- if and
    only if the flag is \textbf{jump-spanned}: every member is the span of the jumps below it. In
    particular this holds in the following two cases:
    \begin{enumerate}
        \item \textbf{Well-ordered flags.} If the flag is well-ordered (no infinite strictly
              descending chain of members), then ``peeling'' terminates: given $v \neq 0$, let
              $W$ be the least member containing it; $W$ is a successor $W^{-} \prec W$ (a limit
              member is the union of its predecessors, so a vector in it already lies in a smaller
              member), so $v = c\, b + w$ with $b$ the jump-vector at $W$ and $w \in W^{-}$, and
              one repeats on $w$. Well-foundedness forbids an infinite descent, so $v$ is
              expressed in finitely many jump-vectors.
        \item \textbf{Cut-flags.} The flags of Proposition~\ref{prop:basis-flag} are jump-spanned
              by construction (each $W_L$ is spanned by the $b_i$ with $i \in L$, i.e.\ by the
              jumps below it), and are exactly the complete flags arising from totally ordered
              bases.
    \end{enumerate}
\end{proposition}

\begin{proof}
    ($\Leftarrow$) First note the flag's greatest member is $V$: the union of all members is a
    subspace (a union of a chain of subspaces) comparable to every member, so by maximality it
    belongs to the chain as its top element, and that top must be $V$ (otherwise $V$ itself could
    be appended). Jump-spannedness applied to this top member gives that $V$ is the span of all
    jump-vectors, so they generate $V$; with Proposition~\ref{prop:jump-independent} they form a
    basis, adapted to the flag since each jump-vector lies in its member.

    ($\Rightarrow$) Suppose the jump-vectors span $V$, and let $W$ be a member. The jump-vectors
    whose jump is $\Subspace[K] W$ clearly lie in $W$. Conversely, take $w \in W$ and write
    $w = \sum_{k \in F} c_k v_k$ over jump-vectors ($F$ finite, every $c_k \neq 0$). Let $v_M$
    have the highest jump $W_M^{-} \prec W_M$ among them. If $W_M \not\Subspace[K] W$, then
    $W \subsetneq W_M$; as $W$ is comparable to $W_M^{-}$ (the members form a chain) and lies
    strictly below $W_M$ with $W_M^{-} \prec W_M$ a covering pair, $W \Subspace[K] W_M^{-}$. Every
    other $v_k$ also lies in $W_M^{-}$ (its jump is lower), so
    $v_M = c_M^{-1}\bigl(w - \sum_{k \in F \setminus \Set{M}} c_k v_k\bigr) \in W_M^{-}$,
    contradicting $v_M \in W_M \setminus W_M^{-}$. Hence $W_M \Subspace[K] W$, every jump in the
    expansion is $\Subspace[K] W$, and $w$ is a combination of jumps below $W$. Thus $W$ is the
    span of the jumps below it. The two displayed cases are verified inline above.
\end{proof}

\begin{remark}[The Jump-Spanned Hypothesis Is Genuine]\label{rem:not-jump-spanned}
    Not every complete flag is jump-spanned, so the equivalence above is not vacuous. A maximal
    chain in $\SubspaceLattice{V}$ need not be well-ordered: it can contain an order-\emph{dense}
    interval --- a subinterval order-isomorphic to a densely ordered family of subspaces, with no
    covering pairs at all. Across such an interval the dimension grows while there are \emph{no}
    jumps, so the jump-vectors cannot reach those members and fail to span $V$; correspondingly
    there is no least member containing a given vector, and ``peeling'' has nothing to start from.
    Thus jump-spannedness --- automatic for cut-flags
    (Proposition~\ref{prop:basis-flag}) and for well-ordered flags --- is a real restriction, and
    recovering a basis from an arbitrary complete flag genuinely requires it.
\end{remark}

\begin{remark}[Round Trip; What Is Indexed by What]\label{rem:roundtrip}
    The passage basis $\rightsquigarrow$ cut-flag $\rightsquigarrow$ jumps returns the original
    basis together with its order $I$. Note that the basis is indexed by the \emph{jumps} of
    $\widehat{I}$ --- which biject with $I$ --- and \emph{not} by all of $\widehat{I}$: the
    gap-filling cuts are continuous (no jump, no vector). So the flag index $\widehat{I}$
    (Dedekind-complete) and the basis index $I$ are genuinely different objects whenever $I$ has
    gaps. Cardinality forces the gap: $\Dim[K]{V} = \Card{I} \leqslant \Card{\widehat{I}}$, e.g.\
    $\Card{\bb{Q}} = \aleph_0$ versus $\Card{\widehat{\bb{Q}}} = \mathfrak{c}$ for the
    $\bb{Q}$-indexed basis of Example~\ref{ex:cut-flag-infinite}.
\end{remark}

\begin{corollary}[Converse in Finite Dimensions]\label{rem:flag-from-basis}
    In finite dimensions every complete flag
    $\{0\} = V_0 \subsetneq V_1 \subsetneq \cdots \subsetneq V_n = V$
    is a finite chain, hence well-ordered and jump-spanned, so by
    Proposition~\ref{prop:jump-recovery} it admits an \emph{adapted ordered basis}
    $(b_0, \ldots, b_{n-1})$ with $V_{i+1} = \Span[K]{\Set{b_0, \ldots, b_i}}$, obtained by
    choosing $b_i \in V_{i+1} \setminus V_i$ at each step. Different choices give different
    adapted bases for the same flag.
\end{corollary}
